\section{The BX3 Three-Layer Model}\label{sec:three-layer}

The \bxthree{} three-layer model is the architectural substrate within which all five
pillars of the framework operate.  It is not a separate mechanism from the pillars ---
it is the structural foundation that makes the pillars possible and that makes their
guarantees unconditional rather than contingent.  This section specifies each layer's
function, the interaction constraints between layers, the data structures maintained
by each layer, and the invariants that the layers collectively enforce.

% ---------------------------------------------------------------
\subsection{Layer Definitions}\label{sec:layer-definitions}

Every \bxthree{} node is organized into three immutable functional layers,
each with a distinct role in the accountability architecture and each with
strictly delineated authority that cannot be exceeded or transferred.

\textbf{The \purpose{} layer} carries intent, judgment, and final authority.
It is the layer at which the human principal's objectives are translated into
operational directives and at which the operating range $R(n)$ for node $n$ is
defined.  The \purpose{} layer receives all escalated exceptions and issues all
binding resolution directives.  It is the exclusive terminus for any exception
state that exceeds machine-actor provisioned authority.  The layer is
non-delegable: no machine actor possesses the judgment capacity to serve as a
substitute terminus for exceptions that exceed its own operating range.
This non-delegability is not a policy constraint -- it is a structural
property of the layer's function.  The \purpose{} layer carries the
principal's intent, and no agent operating under delegated authority can possess
the principal's intent in sufficient completeness to substitute for it.

\textbf{The \bounds{} layer} carries bounded reasoning and constrained execution
within $R(n)$.  It is the layer at which the node evaluates whether proposed
actions fall within its current authority by reference to the authorization
ledger maintained by the \fact{} layer, executes resolution procedures when
exception conditions are detected, and emits the \texttt{bailout\_signal} when
the exception exceeds its current scope.  The \bounds{} layer acts autonomously
within $R(n)$ but is architecturally compelled to escalate when it exits that
range.  The $T_{\text{bounds}}$ timeout parameter is enforced at this layer,
ensuring that bounded reasoning does not become unbounded deliberation.

\textbf{The \fact{} layer} carries deterministic enforcement and forensic
auditability.  It is the layer at which the authorization ledger is consulted,
the state machine transition relations of all five pillars are enforced, and
the append-only Forensic Ledger is maintained.  The \fact{} layer is the
fail-safe anchor of the framework: it is the only layer that may issue a
definitive halt to autonomous action, and it does so by refusing to approve
any proposed action that lacks a valid authorization entry.  No machine actor
may modify, delete, or suppress entries in the Forensic Ledger; the layer's
function is to enforce the record rather than to interpret it.

The three layers are not merely a conceptual taxonomy -- they are enforced
architectural divisions with distinct execution semantics enforced by the
\fact{} layer's pre-commit hook.  A node in the \bounds{} layer cannot issue
authorizations; a node in the \purpose{} layer cannot spawn sub-agents without
a designated human anchor; a node in the \fact{} layer cannot define operating
ranges.  These constraints are not convention -- they are structural invariants
enforced at the pre-commit stage before any layer's code can execute.

% ---------------------------------------------------------------
\subsection{The Authorization Ledger}\label{sec:auth-ledger}

The authorization ledger is the definitive record of all actions that the
\fact{} layer has approved for execution.  Every entry in the ledger is
immutable once written: no machine actor, including the \bounds{} layer or any
parent node, may modify or delete a ledger entry after it is committed.  The
ledger is maintained by the \fact{} layer and is consulted by the \bounds{}
layer before any proposed action is dispatched to the execution engine.

Each ledger entry $e$ is a 4-tuple:
\[
e = (t,\ \text{node}(e),\ \text{action}(e),\ \text{scope}(e))
\]
where $t$ is the logical timestamp, $\text{node}(e)$ is the identifier of
the node that requested the authorization, $\text{action}(e)$ is the action
descriptor, and $\text{scope}(e)$ is the operating range parameter within
which the action was evaluated.  The scope field is critical: it records the
$R(n)$ parameter under which the action was authorized, enabling retrospective
determination of whether the action was taken within or outside the principal's
stated intent at the time of authorization.

The ledger provides the foundation for the \bxthree{} accountability model:
every consequential action is attributable to a node and authorized within a
defined scope, and the complete action history of any node is reconstructable
from the ledger as an immutable chronological record.  The ledger is not a
conventional audit log -- it is an append-only authorization substrate that
the \fact{} layer enforces at the pre-commit stage.

% ---------------------------------------------------------------
\subsection{The Forensic Ledger}\label{sec:forensic-ledger}

The Forensic Ledger is the append-only record of all protocol events across
all five pillars of the framework.  Unlike the authorization ledger, which
records only approved actions, the Forensic Ledger records every state
transition, every trigger condition, every timeout event, every escalation,
every pathway selection decision, and every human directive issued by an
$h(n)$.  The Forensic Ledger is maintained exclusively by the \fact{} layer;
no machine actor may write to, modify, or suppress any entry.

Each Forensic Ledger entry $f$ is a 6-tuple:
\[
f = (\text{id}(f),\ t_{\text{logical}}(f),\ \text{node}(f),\ \text{event}(f),\ \text{context}(f),\ \text{hash}(f))
\]
where $\text{id}(f)$ is the unique identifier of the entry, $t_{\text{logical}}(f)$
is the logical timestamp assigned by the \fact{} layer, $\text{node}(f)$ is the
originating node, $\text{event}(f)$ is the event type, $\text{context}(f)$ is the
full diagnostic context associated with the event (the content of which depends
on the event type), and $\text{hash}(f)$ is the SHA-256 hash of the preceding
entry, forming a cryptographically chained sequence from the genesis entry to
the most recent.  The cryptographic chaining ensures that no entry can be
modified, deleted, or inserted without breaking the chain, a condition detectable
by any auditor without access to the ledger's contents.

The Forensic Ledger serves two distinct functions that are often conflated in
conventional audit systems.  First, it is an \textit{enforcement mechanism}: the
\fact{} layer's pre-commit hooks for the Bailout Protocol, Sandbox Gate, and
Spatial Firewall all query the Forensic Ledger to determine the current state
of their respective state machines before committing any state change.  The
ledger is not merely a record -- it is the authoritative source of truth for
the runtime state of all five pillars.  Second, it is an \textit{evidence
mechanism}: any human principal, regulator, or external auditor can reconstruct
the complete exception and resolution history of any node from the ledger,
including all diagnostic context, all pathway selections, and all human
directives, with cryptographic assurance that the record has not been tampered
with.

The Forensic Ledger is the architectural bridge between the \bxthree{}
Framework's runtime accountability properties and its retrospective auditability
properties.  The Bailout Protocol ensures that every unresolved exception reaches
$h(n)$; the Forensic Ledger ensures that this fact can be demonstrated to any
external auditor at any future point in time.  Together they constitute
\textit{accountable escalation}: the guarantee that exceptions are not merely
resolved but resolved in a way that is permanently attributable.

% ---------------------------------------------------------------
\subsection{Layer Interaction Constraints}\label{sec:layer-constraints}

The three layers interact through strictly defined interfaces that enforce the
separation of concerns essential to the framework's accountability properties.
These interfaces are not conventions -- they are structural constraints
enforced by the \fact{} layer's pre-commit hook.

\begin{enumerate}
\item[\textbf{IC-1.}] The \bounds{} layer may request authorization for an action
by submitting an authorization request to the \fact{} layer.  The \fact{} layer
consults the authorization ledger and returns approval or rejection.  The
\bounds{} layer may not bypass this request-response cycle for any consequential
action.  The \purpose{} layer does not participate in this cycle -- it defines
the scope, it does not execute within it.

\item[\textbf{IC-2.}] The \purpose{} layer receives exceptions from the
\fact{} layer only when the \bounds{} layer has triggered the Bailout Protocol
and the protocol has reached the \purpose{} layer via the escalation algorithm.
The \purpose{} layer does not receive raw exception signals -- it receives
escalated exceptions with full diagnostic context appended by the propagation
chain.

\item[\textbf{IC-3.}] The \fact{} layer writes to the Forensic Ledger for every
state transition of every pillar's state machine and for every event that
requires recording under the pillar's specification.  The \fact{} layer does not
filter, aggregate, or summarize entries -- it records every event in full.
The \purpose{} and \bounds{} layers have no write access to the Forensic Ledger.

\item[\textbf{IC-4.}] The \purpose{} layer issues binding resolution directives
by committing an entry to the authorization ledger through the \fact{} layer.
The directive is not effective until it is recorded in the ledger and verified
by the \fact{} layer's pre-commit hook.  No machine actor may issue, simulate,
or substitute a resolution directive.
\end{enumerate}

These interaction constraints collectively enforce the separation of concerns
that makes the \bxthree{} accountability architecture sound.  The \purpose{}
layer defines intent but does not execute within scope; the \bounds{} layer
executes within scope but cannot authorize actions outside it; the \fact{} layer
enforces both constraints and records both the actions and the exceptions.
No layer can substitute for another, and no layer can perform its functions
without the others.

% ---------------------------------------------------------------
\subsection{Three-Layer Invariants}\label{sec:layer-invariants}

The following invariants hold for every \bxthree{} node throughout its
operational lifetime.  They are established at node provisioning, maintained
by the \fact{} layer's pre-commit hooks, and are inviolable by any machine
actor operating within the framework's architecture.

\begin{invariant}[Intent Preservation]\label{inv:intent}
For every action $a$ executed by node $n$, there exists a ledger entry
$e$ such that $\text{action}(e) = a$ and $\text{scope}(e) \subseteq R(n)$,
where $R(n)$ is the operating range defined by the \purpose{} layer of $n$.
\end{invariant}

\begin{invariant}[Immutable Parent Pointer]\label{inv:parent-pointer}
The parent pointer $\alpha(n)$ of every node $n$ is established at node
provisioning and is never modified during the node's operational lifetime.
No machine actor, including the \bounds{} layer and any parent node, may
modify $\alpha(n)$.
\end{invariant}

\begin{invariant}[Exclusive Human Terminus]\label{inv:human-anchor}
For every node $n$, the human anchor $h(n)$ is established at node
provisioning and is never modified during the node's operational lifetime.
No machine actor may serve as the terminus for an unresolved exception state
at node $n$.  The only permissible terminus is $h(n)$.
\end{invariant}

\begin{invariant}[Authorization Ledger Immutability]\label{inv:ledger-immutability}
Once an entry $e$ is written to the authorization ledger, it is never
modified or deleted.  The ledger is append-only, and its integrity is
enforced by the \fact{} layer's pre-commit hook.
\end{invariant}

\begin{invariant}[Forensic Ledger Completeness]\label{inv:forensic-completeness}
Every state transition, every trigger condition, every timeout event, every
human directive, and every exception event in any pillar's state machine is
recorded in the Forensic Ledger before any subsequent action is taken by any
layer of node $n$.
\end{invariant}

These five invariants are not independent design choices -- they are mutually
defining structural facts.  Invariant~\ref{inv:intent} depends on
Invariant~\ref{inv:ledger-immutability} for its enforcement; Invariant
~\ref{inv:human-anchor} depends on Invariant~\ref{inv:parent-pointer} for the
propagation path; Invariant~\ref{inv:forensic-completeness} depends on
Invariant~\ref{inv:ledger-immutability} for the integrity of the recorded
context.  Together they constitute the structural foundation for the upstream
accountability guarantee: every consequential action is within scope, every
exception is attributable, and every unresolved exception reaches $h(n)$.
